This section formalises rock--TBM interpretation as an excavation-step-indexed spatial entity--relation graph. The method starts from a TSP-derived rock voxel field, an advancing parameterized TBM surface and aligned monitoring records under a common excavation-step index. It then constructs a sparse dynamic rock--TBM spatial relation graph at each excavation step and aggregates edge-weighted geological anomaly scores into a component-level geological anomaly index. The graph is sparse because candidate rock--TBM relations are retained only when the four spatial and excavation-state constraints are jointly satisfied.

Given geological voxels $D_{geo}$, a parameterized TBM surface $M_{TBM}$ and excavation step $t$, the graph-construction operator is
\begin{equation}
G : (D_{geo}, M_{TBM}, t) \mapsto G_t .
\end{equation}
The field-computation operator is
\begin{equation}
F : (G_t, \phi(g_i), \rho_j) \mapsto \{I_c(t), \Delta I_c(t)\}_{c\in C,t=1,\ldots,T},
\end{equation}
where $\phi(g_i)$ maps geological attributes to a rock-voxel anomaly score and $\rho_j$ is the TBM component label. Monitoring records are used to read the resulting index against excavation responses; they are not inputs to $G$ and do not define graph edges.


\subsection{Rock, machine and monitoring data representation}
The input data are organised into two graph entity types and aligned monitoring records: rock voxels, TBM surface nodes and excavation monitoring responses. Monitoring records are not graph nodes; they are used after graph computation to evaluate and interpret the component-indexed geological anomaly index.

\textbf{Rock voxel.} The geological input is a TSP-derived voxel field $D_{geo}=\{(c_i,g_i)\}$, where $c_i=(x_i,y_i,z_i)$ is the coordinate of rock voxel $i$ and $g_i$ is its geological attribute vector, such as P-wave velocity $V_p$, S-wave velocity $V_s$, dynamic Young's modulus, velocity ratio, Poisson's ratio or density. Rock voxels are the geological spatial support of the graph.

\textbf{TBM surface node.} The machine input is a parameterized TBM surface. Each surface node is represented by $(p_j(0),n_j,\rho_j)$, where $p_j(0)$ is the initial surface-node coordinate, $n_j$ is the local outward surface normal and $\rho_j$ is the machine component label. The current implementation partitions the machine surface into cutterhead, front shield, middle shield and tail shield components.

\textbf{Aligned monitoring records.} At each excavation step $t$, the monitoring vector is
\begin{equation}
u_t=[\mathrm{AdvanceRate}_t,\mathrm{RPM}_t,\mathrm{Torque}_t,\mathrm{Thrust}_t,\mathrm{Penetration}_t,\mathrm{ShieldPressure}_t].
\end{equation}
The response variables used downstream are denoted $r_t^{(k)}$, where $k$ indexes a response channel. Monitoring responses are aligned to $I_c(t)$ and $\Delta I_c(t)$; they do not participate in graph construction.

\textbf{Common chainage coordinate.} TSP coordinates, TBM geometry and monitoring records are mapped to a common local excavation coordinate. With $\mathrm{ShieldMileage}(t)$ the shield mileage stamp at step $t$,
\begin{equation}
x_{local}(t)=\mathrm{ShieldMileage}(t)-\mathrm{ShieldMileage}_{start},
\end{equation}
and TBM surface nodes are advanced along the excavation direction $d$:
\begin{equation}
p_j(t+1)=p_j(t)+\Delta x d.
\end{equation}

\begin{table}[t]
\centering
\caption{Entities, relations and indices in the revised rock--TBM spatial relation graph.}
\label{tab:entity_model}
\begin{tabular}{lll}
\toprule
Entity or quantity & Definition & Role \\
\midrule
$V_t^r$ & Active TSP-derived rock voxels & Geological spatial support \\
$V_t^m$ & Component-labelled TBM surface nodes & Machine-side spatial support \\
$E_t^{rr}$ & Rock--rock voxel neighbourhoods & Geological contiguity context \\
$E_t^{mm}$ & TBM surface neighbourhoods & Machine-surface topology context \\
$E_t^{rm}$ & Screened rock--TBM candidate relations & Candidate relation support \\
$A_c(t)$ & Sum of geometry weights for component $c$ & Geometric exposure metric \\
$I_c(t)$ & Edge-weighted anomaly intensity for component $c$ & Main component-level index \\
$e_{t+h}^{(k)}$ & Persistence residual of response $k$ & Aligned monitoring residual \\
\bottomrule
\end{tabular}
\end{table}

\begin{figure}[t]
\centering
\includegraphics[width=0.95\linewidth]{fig_method_workflow.png}
\caption{Workflow of the dynamic rock--TBM spatial relation graph. Rock voxels, TBM surface nodes and aligned monitoring records are organised under a common excavation-step coordinate. Candidate rock--TBM relations are generated under active-zone, distance, surface-normal compatibility and excavation-state constraints, aggregated into component-level geological anomaly indices, and queried back to edge-level provenance.}
\label{fig:overview}
\end{figure}

\subsection{Geometry-constrained sparse graph construction}
At each excavation step, a graph snapshot is defined as
\begin{equation}
G_t=(V_t^r\cup V_t^m,E_t^{rr}\cup E_t^{mm}\cup E_t^{rm}).
\end{equation}
Rock node features are
\begin{equation}
x_i^r(t)=[c_i,g_i,s_i(t)],
\end{equation}
where $s_i(t)$ records whether voxel $i$ remains active and eligible for relation construction at step $t$. TBM node features are
\begin{equation}
x_j^m(t)=[p_j(t),n_j(t),\rho_j].
\end{equation}
The intra-rock edge set $E_t^{rr}$ is constructed from voxel neighbourhoods, and $E_t^{mm}$ is constructed from local neighbourhoods on the parameterized TBM surface. The central design decision is the construction of $E_t^{rm}$, the rock--TBM candidate relation set.

A candidate edge $(i,j)$ is included in $E_t^{rm}$ only when four conditions are jointly satisfied:
\begin{equation}
\begin{cases}
c_i\in \Omega_t,\\
d_{ij}(t)\leq \tau_{edge},\\
\kappa_{ij}(t)\geq \eta_{min},\\
s_i(t)=1.
\end{cases}
\end{equation}
The first rule restricts relations to the excavation-relevant active rock volume $\Omega_t$; the second uses
\begin{equation}
d_{ij}(t)=\|c_i-p_j(t)\|
\end{equation}
to retain spatially close rock--machine pairs; the third uses the signed surface-normal compatibility
\begin{equation}
\kappa_{ij}(t)=\max\left(0,\frac{n_j(t)^\top(c_i-p_j(t))}{d_{ij}(t)+\epsilon}\right)
\end{equation}
to retain geometrically plausible facing relations; and the fourth prevents already excavated voxels from repeatedly contributing to later graph snapshots.

The active zone separates face-contact plausibility from shield-side proximity:
\begin{equation}
\Omega_t=\Omega_t^{face}\cup\Omega_t^{shield},
\end{equation}
with
\begin{equation}
\Omega_t^{face}=\{i\mid 0\leq x_i-x_f(t)\leq L_f,\; r_i\leq R_c+\tau_{zone}\},
\end{equation}
\begin{equation}
\Omega_t^{shield}=\{i\mid x_{tail}(t)\leq x_i\leq x_f(t),\; R_s\leq r_i\leq R_s+\tau_{zone}\}.
\end{equation}
For each retained candidate relation, the geometry-compatibility weight is
\begin{equation}
w_{ij}^{rm}(t)=\exp\left(-\frac{d_{ij}(t)}{\tau_{edge}}\right)\kappa_{ij}(t).
\end{equation}
This quantity is a spatial relation weight; it is not a contact force, pressure or friction estimate.

\begin{figure}[t]
\centering
\fbox{\parbox{0.92\textwidth}{\centering Placeholder for candidate-relation construction: active-zone, distance, normal-compatibility and excavation-state filtering.}}
\caption{Geometry-constrained rock--TBM candidate relation construction. The graph is obtained by screening all possible rock voxel--TBM surface-node pairs through active-zone membership, distance, surface-normal compatibility and excavation-state constraints.}
\label{fig:construction}
\end{figure}

\subsection{Component-level geological anomaly index}
The geological contribution of each rock voxel is represented by a bounded TSP-derived anomaly score
\begin{equation}
q_i=\phi(g_i).
\end{equation}
In the main formulation, $q_i$ is a low-velocity anomaly score derived from the training active-zone distribution of P-wave velocity:
\begin{equation}
q_i=\mathrm{clip}\left(\frac{Q_{95}^{train}(V_p)-V_{p,i}}{Q_{95}^{train}(V_p)-Q_{5}^{train}(V_p)+\epsilon},0,1\right).
\end{equation}
This fixed training-reference scaling avoids per-chainage min--max normalisation and keeps values comparable along chainage; lower $V_p$ values correspond to higher anomaly scores under the common reference distribution.

For each TBM component $c$, let $V_c^m$ denote the TBM surface nodes with component label $c$. The geometric exposure is
\begin{equation}
A_c(t)=\sum_{j\in V_c^m}\sum_{i:(i,j)\in E_t^{rm}} w_{ij}^{rm}(t),
\end{equation}
and the component-level geological anomaly index is
\begin{equation}
I_c(t)=\frac{\sum_{j\in V_c^m}\sum_{i:(i,j)\in E_t^{rm}}w_{ij}^{rm}(t)q_i}{A_c(t)+\epsilon}.
\end{equation}
The index is computed for the cutterhead, front shield, middle shield and tail shield at each excavation step.

Because the confirmed case is shield jamming, shield components are also summarised as an exposure-weighted shield-group index:
\begin{equation}
I_{shield}^{A}(t)=\frac{\sum_{c\in C_{shield}}A_c(t)I_c(t)}{\sum_{c\in C_{shield}}A_c(t)+\epsilon},
\quad C_{shield}=\{\mathrm{front},\mathrm{middle},\mathrm{tail}\}.
\end{equation}
This replaces the unweighted arithmetic mean and avoids assigning equal influence to components with different relation supports. A shield-specificity contrast is further defined as
\begin{equation}
S_{shield}(t)=\frac{\sum_{c\in C_{shield}}A_c(t)I_c(t)}{A_{cutterhead}(t)I_{cutterhead}(t)+\epsilon}.
\end{equation}
The step-level change is
\begin{equation}
\Delta I_c(t)=I_c(t)-I_c(t-1).
\end{equation}

For response alignment, short-horizon TBM responses are evaluated through persistence residuals rather than raw responses:
\begin{equation}
e_{t+h}^{(k)}=r_{t+h}^{(k)}-r_t^{(k)}.
\end{equation}
Association statistics between $I_c(t)$ and $e_{t+h}^{(k)}$ are interpreted as exploratory index--residual co-variation within the evaluated chainage interval, not as causal attribution.

For edge-level provenance, candidate edges for any diagnostic sample are ranked by
\begin{equation}
s_{ij}(t)=w_{ij}^{rm}(t)q_i.
\end{equation}
Each record retains rock voxel identity, TBM surface-node identity, component label $\rho_j$, distance $d_{ij}(t)$, normal compatibility $\kappa_{ij}(t)$, geometry weight $w_{ij}^{rm}(t)$, anomaly score $q_i$ and weighted contribution $s_{ij}(t)$.
